
If \(S\) is a prescheme, \(H\) a commutative \(S\)-group
prescheme, and \(G\) an \(S\)-group prescheme acting on \(H\), we
denote by \(H^i(G,H)\) the \(i\)-th Hochschild cohomology group of \(G\)
with values in \(H\) (Appendix I, 1).

To prove \SGARefLink{sga3:xvii:theorem:5-1-1}{5.1.1}, we shall proceed in several steps.

\SGASubsection{5.2}{\texorpdfstring{Proof of \SGARefLink{sga3:xvii:theorem:5:1:1:item:i}{5.1.1(i)} and \SGARefLink{sga3:xvii:theorem:5:1:1:item:ii}{(ii)}
when \(U\) is smooth and \(H\) is étale}{Proof of \SGARefLink{sga3:xvii:theorem:5:1:1:item:i}{5.1.1(i)} and
\SGARefLink{sga3:xvii:theorem:5:1:1:item:ii}{(ii)} when U is smooth and H is etale}}
\SGARefTarget{sga3:xvii:subsection:5:2}

\SGAStatement{Lemma 5.2.1}{With the notation of \SGARefLink{sga3:xvii:theorem:5-1-1}{5.1.1}, suppose
that \(H\) is étale.  The canonical morphism \(E\to H\) has a
section \(s:H\to E\), defined over \(k\), in each of the following
cases:
\begin{enumerate}
\item[(a)] \SGARefTarget{sga3:xvii:lemma:5:2:1:item:a}
  \(k\) is algebraically closed;
\item[(b)] \SGARefTarget{sga3:xvii:lemma:5:2:1:item:b}
  \(k\) is perfect and \(U\) is smooth and connected;
\item[(c)] \SGARefTarget{sga3:xvii:lemma:5:2:1:item:c}
  \(U\) is \(k\)-solvable.
\end{enumerate}}
\SGARefTarget{sga3:xvii:lemma:5:2:1}

Assertion \SGARefLink{sga3:xvii:appendix:A}{(a)} follows simply from the surjectivity of
\(E(k)\to H(k)\).  We have \SGARefLink{sga3:xvii:section:b}{(b)}\(\Rightarrow\)\SGARefLink{sga3:xvii:appendix:C}{(c)} by \SGARefLink{sga3:xvii:lemma:4:1:2:item:b}{4.1.2(b)}
and \SGARefLink{sga3:xvii:lemma:2:3:bis}{2.3 bis}, so it is enough to treat case \SGARefLink{sga3:xvii:appendix:C}{(c)}.  Let \(x\) be a
point of \(H\).  It is an open and closed subset of \(H\), and
the induced subscheme is isomorphic to \(\operatorname{Spec}K\),
where \(K\) is a finite separable extension of \(k\).  Let \(X\)
be the inverse image in \(E\) of the scheme \(x\).  The
\(K\)-scheme \(X\) is a principal homogeneous \(K\)-space under
\(U_K\).  But \(U_K\) has a composition series whose successive
quotients are isomorphic to \((\mathbb G_a)_K\); hence
\[
  H^1(K,U_K)=0
\]
(J.-P. Serre, \emph{Cohomologie galoisienne}, Chapter III,
\SGARefLink{sga3:xvii:section:6}{Proposition 6}).  Consequently \(X\) is trivial and has a
\(K\)-rational point.  Thus, over every point \(x\) of \(H\), we
obtain a \(k\)-section of \(E\to H\), and these give a section
\(H\to E\).

\SGAStatement{\SGAInlineRefTarget{sga3:xvii:lemma:5-2-3}Lemma 5.2.3}{With the notation of \SGARefLink{sga3:xvii:theorem:5-1-1}{5.1.1}, suppose that
\(H\) is étale.  Then:
\begin{enumerate}
\item[(i)] \SGARefTarget{sga3:xvii:lemma:5:2:3:item:i}
  The extension \(E\) is trivial in each of the following cases:
  \begin{enumerate}
  \item[(a)] \SGARefTarget{sga3:xvii:lemma:5:2:3:item:i:a}
    \(U\) is commutative and \(E\to H\) has a section;
  \item[(b)] \SGARefTarget{sga3:xvii:lemma:5:2:3:item:i:b}
    \(k\) is algebraically closed;
  \item[(c)] \SGARefTarget{sga3:xvii:lemma:5:2:3:item:i:c}
    \(k\) is perfect and \(U\) is connected;
  \item[(d)] \SGARefTarget{sga3:xvii:lemma:5:2:3:item:i:d}
    \(U\) is \(k\)-solvable.
  \end{enumerate}
  Moreover, in each of the preceding cases, two lifts \(H'\) and
  \(H''\) of \(H\) in \(E\) are conjugate by an element of
  \(U(k)\).
\item[(ii)] \SGARefTarget{sga3:xvii:lemma:5:2:3:item:ii}
  Let \(R\) be the \(k\)-functor
  \[
    (\operatorname{Sch}/k)^\circ\longrightarrow\operatorname{Ens}
  \]
  for which, for every \(k\)-prescheme \(S\), \(R(S)\) is the set
  of lifts of \(H_S\) in \(E_S\).  If \(U\) is commutative, then
  \(R\) is representable by a nonempty affine \(k\)-scheme.  The
  group \(U\) acts on \(R\) by inner automorphisms, and this action
  makes \(R\) a homogeneous space under \(U\) for the fpqc
  topology (cf.\ Exposé IV).
\end{enumerate}}
\SGARefTarget{sga3:xvii:lemma:5:2:3}

\SGAProof{Proof of \SGARefLink{sga3:xvii:lemma:5:2:3:item:i}{(i)}}
We shall reduce cases \SGARefLink{sga3:xvii:section:b}{(b)}, \SGARefLink{sga3:xvii:appendix:C}{(c)}, and \SGARefLink{sga3:xvii:lemma:5:2:3:item:i:d}{(d)} to case \SGARefLink{sga3:xvii:appendix:A}{(a)}.

\medskip\noindent\textit{Case \SGARefLink{sga3:xvii:section:b}{(b)}.}\quad
Since \(U\) has a characteristic composition series with
commutative successive quotients
(\SGARefLink{sga3:xvii:theorem:3-5}{3.5}, \SGARefLink{sga3:xvii:theorem:3:5:item:i}{(i)}\(\Leftrightarrow\)\SGARefLink{sga3:xvii:theorem:3:5:item:iv}{(iv)}), we immediately reduce to the
case in which \(U\) is commutative.  Moreover, \(E\to H\) has a
section by \SGARefLink{sga3:xvii:lemma:5:2:1:item:a}{5.2.1(a)}, so we are reduced to case \SGARefLink{sga3:xvii:appendix:A}{(a)}.

\medskip\noindent\textit{Case \SGARefLink{sga3:xvii:lemma:5:2:3:item:i:d}{(d)}.}\quad
We argue in the same way, using \SGARefLink{sga3:xvii:definition:5:1:0:item:ii}{5.1.0(ii)} and \SGARefLink{sga3:xvii:lemma:5:2:1:item:c}{5.2.1(c)}.

\medskip\noindent\textit{Case \SGARefLink{sga3:xvii:appendix:C}{(c)}.}\quad
Let \(E_{\mathrm{red}}\) be the reduced subgroup associated with
\(E\).  Since \(H\) is smooth, \(E_{\mathrm{red}}\) is an
extension of \(H\) by
\[
  U_0=U\cap E_{\mathrm{red}},
\]
and we may replace \(E\) by \(E_{\mathrm{red}}\).  The group \(U\),
and therefore \(U_0\), is connected, and \(U_0\) is plainly the
identity component of \(E_{\mathrm{red}}\); in particular it is
smooth.  Thus \(U_0\) is smooth and connected and \(k\) is
perfect, so \(U_0\) is \(k\)-solvable by \SGARefLink{sga3:xvii:remark:4:1:4:item:b}{4.1.4(b)}, reducing us to
case \SGARefLink{sga3:xvii:lemma:5:2:3:item:i:d}{(d)}.

The preceding reductions show that in cases \SGARefLink{sga3:xvii:section:b}{(b)}, \SGARefLink{sga3:xvii:appendix:C}{(c)}, and \SGARefLink{sga3:xvii:lemma:5:2:3:item:i:d}{(d)} we
may suppose that \(U\) has a characteristic composition series
\((U_i)\) such that \(U_i/U_{i+1}\) is commutative and the maps
\[
  U_i(k)\longrightarrow (U_i/U_{i+1})(k)
\]
are surjective.  (In cases \SGARefLink{sga3:xvii:appendix:C}{(c)} and \SGARefLink{sga3:xvii:lemma:5:2:3:item:i:d}{(d)}, the latter assertion
follows from \(H^1(k,\mathbb G_a)=0\).)  An immediate dévissage
then shows that it is enough to prove the conjugacy of two lifts
\(H'\) and \(H''\) in \(E\) when \(U\) is commutative.  In short,
it is enough to prove \SGARefLink{sga3:xvii:lemma:5:2:3:item:i}{(i)}\SGARefLink{sga3:xvii:lemma:5:2:4:item:a}{(a)}.  In that case the extension \(E\)
is trivial provided that \(H^2(H,U)=0\) (Appendix I, \SGARefLink{sga3:i:definition:3-1}{3.1}), and
\(H'\) and \(H''\) are conjugate provided that \(H^1(H,U)=0\).
We now use the following lemma.

\SGAStatement{Lemma 5.2.4}{Let \(S\) be a prescheme, \(U\) a
commutative \(S\)-group prescheme, and \(H\) an étale finite
\(S\)-group prescheme of rank \(n\) acting on \(U\).  Then the
groups \(H^i(H,U)\), for \(i>0\), are annihilated by \(n\) in
each of the following cases:
\begin{enumerate}
\item[(a)] \SGARefTarget{sga3:xvii:lemma:5:2:4:item:a}
  \(H\) is a constant \(S\)-group (\SGARefLink{sga3:i:subsection:4-1}{Exposé I, 4.1});
\item[(b)] \SGARefTarget{sga3:xvii:lemma:5:2:4:item:b}
  \(S\) is the spectrum of a field.
\end{enumerate}}
\SGARefTarget{sga3:xvii:lemma:5:2:4}

\SGAProof{Proof of \SGARefLink{sga3:xvii:appendix:A}{(a)}}
By hypothesis, \(H\) is the constant group associated with an
ordinary group \(\{H\}\) of order \(n\).  It is
clear that \(H^i(H,U)\) is isomorphic
to the \(i\)-th cohomology group
\[
  H^i(\{H\},U(S))
\]
of the group \(\{H\}\) with values in the ordinary group \(U(S)\).
It is classical that these groups are annihilated by \(n\)
(J.-P. Serre, \emph{Corps locaux}, Chapter VIII, \SGARefLink{sga3:viii:section:4}{Proposition 4},
Corollary 1).

\SGAProof{Proof of \SGARefLink{sga3:xvii:section:b}{(b)}}
Let \(x\in H^i(H,U)\), where \(i>0\), and represent it by an
\(i\)-cocycle
\[
  f:H^{(i)}\longrightarrow U,
\]
where \(H^{(i)}\) denotes the product over \(k\) of \(i\) copies
of \(H\).  If \(K\) is a finite Galois extension of \(k\) which
splits \(H\), assertion \SGARefLink{sga3:xvii:lemma:5:2:4:item:a}{(a)} shows that \(nf_K\) is a coboundary.
More precisely, a direct computation shows that, if one defines
the \(K\)-morphism
\[
  F_K:H_K^{(i-1)}\longrightarrow U_K
\]
by
\[
  F_K(h_1,\ldots,h_{i-1})
    =\sum_{h\in H(K)}
      f_K(h_1,\ldots,h_{i-1},h),
\]
then, up to sign,
\[
  d(F_K)=nf_K,
  \qquad d\text{ denoting the coboundary operator}.
\]
An immediate Galois argument shows that \(F_K\) descends to a
\(k\)-morphism \(F:H^{(i-1)}\to U\), and consequently
\[
  d(F)=nf.
\]
